% === UNIT 1 ===
\chapter{Unit 1: Exploring One-Variable Data}

\section{SECTION A - Unit Overview}
\begin{lessonbox}[Unit 1 Vocabulary]
\begin{itemize}
    \item \textbf{Categorical Variable}: A variable that places an individual into one of several groups or categories.
    \item \textbf{Quantitative Variable}: A variable that takes numerical values for which it makes sense to find an average.
    \item \textbf{Distribution}: Tells us what values a variable takes and how often it takes these values.
    \item \textbf{Outlier}: An individual value that falls outside the overall pattern of a distribution.
    \item \textbf{Symmetric}: A distribution is symmetric if the right and left sides of the graph are approximately mirror images of each other.
    \item \textbf{Skewed}: A distribution is skewed to the right if the right tail is much longer than the left tail, and skewed to the left if the left tail is much longer.
    \item \textbf{Mean}: The arithmetic average of a distribution, obtained by adding the observations and dividing by the number of observations ($\bar{x}$).
    \item \textbf{Median}: The midpoint of a distribution, the number such that half the observations are smaller and the other half are larger.
    \item \textbf{Standard Deviation}: A measure of how much the values in a distribution typically vary from the mean ($s_x$).
    \item \textbf{Variance}: The square of the standard deviation ($s_x^2$).
    \item \textbf{Z-score}: Tells us how many standard deviations a value is from the mean.
    \item \textbf{Percentile}: The percent of values in a distribution that are less than the individual's data value.
\end{itemize}
\end{lessonbox}

\textbf{AP Exam Weight}: 15-23\% of the AP Statistics Exam.

\textbf{What you will learn}: In this unit, you will learn how to analyze and summarize one-variable data using graphs and numerical summaries. You'll master the distinction between categorical and quantitative data, construct appropriate visual representations, calculate measures of center and spread, and use the normal distribution to determine probabilities.

\section{SECTION B - Core Concept Lessons}

\begin{lessonbox}[Lesson 1.1: Categorical vs. Quantitative Data]
\textbf{Definition}: Data represents information collected on variables. A \textbf{categorical variable} places individuals into groups (e.g., eye color, zip code). A \textbf{quantitative variable} takes numerical values for which calculating an average makes sense (e.g., height, age).

\textbf{Context}: Imagine surveying a high school class. Asking students for their favorite subject produces categorical data. Asking for their GPA produces quantitative data. Note that some numbers act as categories—a zip code or a student ID number doesn't measure quantity, so it's categorical!
\end{lessonbox}

\begin{lessonbox}[Lesson 1.2: Measures of Center (Mean vs. Median)]
\textbf{Definition \& Formula}:
The \textbf{mean} ($\bar{x}$) is the arithmetic average: $\bar{x} = \frac{\sum x_i}{n}$.
The \textbf{median} is the middle value when data is ordered.

\textbf{Explanation}: The mean is sensitive to extreme values (non-resistant), while the median is resistant to extreme values. In a skewed right distribution, the mean is pulled toward the right tail, making it larger than the median.

\textbf{Context}: If you look at the starting salaries of 10 teachers and 1 billionaire CEO, the mean salary will be highly inflated by the CEO, making it an inaccurate representation of a "typical" teacher. The median salary would be much more representative.
\end{lessonbox}

\begin{lessonbox}[Lesson 1.3: Measures of Spread (IQR and Standard Deviation)]
\textbf{Definition \& Formula}:
\textbf{Interquartile Range (IQR)} measures the spread of the middle 50\% of data: $IQR = Q_3 - Q_1$.
\textbf{Standard Deviation ($s_x$)} measures typical distance from the mean: $s_x = \sqrt{\frac{\sum (x_i - \bar{x})^2}{n-1}}$.

\textbf{Explanation}: Just like the median, the IQR is resistant to outliers. Standard deviation, however, is heavily influenced by outliers because it depends on the mean.

\textbf{Context}: A pizza delivery shop promises times centered at 25 minutes. If standard deviation is 2 minutes, times are consistently close to 25. If standard deviation is 15 minutes, you might get cold pizza in 40 minutes!
\end{lessonbox}

\begin{lessonbox}[Lesson 1.4: Z-scores and Normal Distributions]
\textbf{Definition \& Formula}: A \textbf{z-score} normalizes data by measuring the number of standard deviations a value $x$ is from the mean $\mu$: $z = \frac{x - \mu}{\sigma}$.

\textbf{Explanation}: Normal distributions are symmetric, bell-shaped curves defined by $\mu$ and $\sigma$. The Empirical Rule states that 68\%, 95\%, and 99.7\% of data fall within 1, 2, and 3 standard deviations of the mean, respectively.

\textbf{Context}: To compare an SAT score (out of 1600) and an ACT score (out of 36), calculate z-scores for both. If your SAT z-score is 1.5 and ACT z-score is 2.0, you performed relatively better on the ACT!
\end{lessonbox}

\section{SECTION C - Worked Examples}

\textbf{Example 1: Identifying Outliers}
\textit{Context}: The number of home runs hit by 10 baseball players are: 5, 8, 12, 14, 15, 18, 22, 24, 25, 45. Find if 45 is an outlier.
\textit{Solution}:
1. Find $Q_1$ and $Q_3$: Median is $(15+18)/2 = 16.5$. $Q_1 = 12$, $Q_3 = 24$.
2. Calculate IQR: $IQR = 24 - 12 = 12$.
3. Check upper boundary: $Q_3 + 1.5 \times IQR = 24 + 1.5(12) = 24 + 18 = 42$.
4. Conclusion: Since $45 > 42$, 45 is indeed an outlier.

\textbf{Example 2: Comparing Mean and Median}
\textit{Context}: A neighborhood has home prices (in thousands): \$250, \$260, \$275, \$300, \$1,500. Calculate the mean and median.
\textit{Solution}:
1. Median = middle value = \$275,000.
2. Mean = $\frac{250+260+275+300+1500}{5} = \frac{2585}{5} = \$517,000$.
3. Because the distribution is right-skewed (by the \$1.5M mansion), the mean is significantly larger than the median.

\textbf{Example 3: Calculating Z-Scores}
\textit{Context}: Average high school female height is $\mu = 64$ inches with $\sigma = 2.5$ inches. Maria is 69 inches tall.
\textit{Solution}:
1. Formula: $z = \frac{x - \mu}{\sigma}$
2. Substitute: $z = \frac{69 - 64}{2.5} = \frac{5}{2.5} = 2.0$
3. Conclusion: Maria is 2 standard deviations above the mean height.

\textbf{Example 4: Using the Empirical Rule}
\textit{Context}: The lifespan of a specific battery is normally distributed with $\mu = 40$ hours and $\sigma = 5$ hours. What percent of batteries last between 30 and 45 hours?
\textit{Solution}:
1. Convert 30 to z-score: $z = (30-40)/5 = -2$.
2. Convert 45 to z-score: $z = (45-40)/5 = 1$.
3. Area between $z=-2$ and $z=0$ is $95\% / 2 = 47.5\%$.
4. Area between $z=0$ and $z=1$ is $68\% / 2 = 34\%$.
5. Total percent: $47.5\% + 34\% = 81.5\%$.

\section{SECTION D - TI-84 CE Keystroke Playbook}
\begin{calculatorbox}[TI-84 Procedures for Unit 1]
\textbf{1. 1-Var Stats (Mean, Median, Standard Deviation)}
\begin{itemize}
    \item Enter data: \texttt{STAT} $\rightarrow$ \texttt{1:Edit} $\rightarrow$ Enter data in L1.
    \item Calculate: \texttt{STAT} $\rightarrow$ \texttt{CALC} $\rightarrow$ \texttt{1:1-Var Stats}.
    \item Ensure List is L1, FreqList is clear, then \texttt{Calculate}.
\end{itemize}

\textbf{2. Creating a Boxplot}
\begin{itemize}
    \item Enter data in L1.
    \item Press \texttt{2ND} $\rightarrow$ \texttt{Y=} (STAT PLOT) $\rightarrow$ \texttt{Plot1}.
    \item Turn Plot ON. For Type, select the modified boxplot (4th icon, shows outliers).
    \item Press \texttt{ZOOM} $\rightarrow$ \texttt{9:ZoomStat} to view.
\end{itemize}

\textbf{3. normalcdf (Finding Area/Probability under Normal Curve)}
\begin{itemize}
    \item \texttt{2ND} $\rightarrow$ \texttt{VARS} (DISTR) $\rightarrow$ \texttt{2:normalcdf(}
    \item Syntax: \texttt{normalcdf(lower, upper, $\mu$, $\sigma$)}
    \item Use -1E99 for a left tail of negative infinity, and 1E99 for a right tail.
\end{itemize}
\end{calculatorbox}

\section{SECTION E - Score-4 FRQ Templates}

\begin{frqrubricbox}[FRQ Template 1: Describing a Distribution]
\textbf{Prompt Type:} "Describe the distribution of..."

\textbf{Template:}
"The distribution of [context variable] is [shape] with a single peak/multiple peaks. The center, as measured by the median, is approximately [value + units]. The spread of the data ranges from [minimum] to [maximum], with an IQR of [IQR value]. There [appears to be / does not appear to be] an outlier at [value + units]."

\textbf{Grader Rubric (4 points):}
\begin{itemize}
    \item \textbf{S} (Shape): Must explicitly state if skewed or symmetric.
    \item \textbf{O} (Outliers): Must mention presence or absence of outliers.
    \item \textbf{C} (Center): Must provide a measure of center (median/mean) in context.
    \item \textbf{S} (Spread): Must describe variability (range/IQR/std dev) in context.
    \item \textit{Note: Context must be woven throughout the response to earn full credit.}
\end{itemize}
\end{frqrubricbox}

\begin{frqrubricbox}[FRQ Template 2: Comparing Two Distributions]
\textbf{Prompt Type:} "Compare the distributions of A and B..."

\textbf{Template:}
"The distribution of [Group A context] is [Shape A], whereas the distribution of [Group B context] is [Shape B]. The center of [Group A] is [greater than/less than/similar to] [Group B], with a median of [Value A] compared to [Value B]. The spread of [Group A] is [greater than/less than] [Group B], with a range of [Range A] vs [Range B]. [Group A] has an outlier at [Value], while [Group B] has no outliers."

\textbf{Grader Rubric (4 points):}
\begin{itemize}
    \item Must explicitly use comparison words ("greater than", "less than", "similar to"). Simply listing the SOCS for both groups without comparing them will max out at a score of 2!
    \item Must compare Centers.
    \item Must compare Spreads.
    \item Must compare Shapes/Outliers in context.
\end{itemize}
\end{frqrubricbox}

\section{SECTION F - Practice Problem Set}
\begin{workbookbox}[Unit 1 Practice Problems]
\textbf{Multiple Choice}
1. Which of the following is a quantitative variable?
(A) Zip code (B) Race (C) Weight (D) Eye color
\textit{Solution: (C). Weight is numerical and an average makes sense.}

2. In a strongly left-skewed distribution, which is true?
(A) Mean > Median (B) Mean < Median (C) Mean = Median (D) Mean is zero
\textit{Solution: (B). The tail on the left pulls the mean down.}

3. If every value in a dataset is increased by 5, the standard deviation:
(A) Increases by 5 (B) Decreases by 5 (C) Remains unchanged (D) Multiplies by 5
\textit{Solution: (C). Adding a constant shifts the center but does not affect spread.}

4. Which measure is most resistant to outliers?
(A) Mean (B) Standard Deviation (C) Range (D) IQR
\textit{Solution: (D). IQR only looks at the middle 50\%, ignoring extremes.}

5. Data set: 10, 12, 14, 15, 18, 22, 28. What is the median?
(A) 14 (B) 15 (C) 18 (D) 17.5
\textit{Solution: (B). Middle value of the ordered 7 numbers is 15.}

6. A z-score of -1.5 means the data point is:
(A) 1.5 above mean (B) 1.5 below mean (C) Invalid (D) The mean
\textit{Solution: (B).}

7. Using empirical rule, what percent falls within 2 std dev?
(A) 68\% (B) 95\% (C) 99.7\% (D) 50\%
\textit{Solution: (B).}

8. If Variance is 16, what is Standard Deviation?
(A) 256 (B) 8 (C) 4 (D) 2
\textit{Solution: (C). $\sqrt{16} = 4$.}

9. Which graphical display is best for checking normality?
(A) Pie chart (B) Bar chart (C) Normal probability plot (D) Scatterplot
\textit{Solution: (C).}

10. A distribution with one peak is called:
(A) Bimodal (B) Unimodal (C) Uniform (D) Skewed
\textit{Solution: (B).}

\textbf{Free Response}
11. The test scores of 10 students are: 45, 78, 82, 85, 88, 90, 91, 93, 95, 99. Calculate the five-number summary and determine if 45 is an outlier.
\textit{Solution: Min=45, Q1=82, Med=89, Q3=93, Max=99. IQR = 93-82 = 11. Lower bound = $82 - 1.5(11) = 65.5$. Since $45 < 65.5$, 45 is an outlier.}

12. A normal distribution has $\mu=50, \sigma=4$. Find the z-score for $x=58$.
\textit{Solution: $z = \frac{58-50}{4} = 2$.}

13. Explain why we use standard deviation instead of variance for describing spread.
\textit{Solution: Standard deviation is in the same units as the original data, making it interpretable, whereas variance is in squared units.}

14. Sketch a density curve that is right-skewed and label the relative positions of the mean and median.
\textit{Solution: [Sketch description] A curve with a long tail on the right. The peak represents the mode, the median is to the right of the peak, and the mean is furthest to the right, pulled by the tail.}

15. If a dataset has $\mu=100, \sigma=15$, what value represents the 97.5th percentile?
\textit{Solution: 97.5th percentile is $z=2$. $x = \mu + z\sigma = 100 + 2(15) = 130$.}
\end{workbookbox}

\section{SECTION G - Unit Summary}
\begin{trapbox}[Common Mistakes]
- Do not mix up categorical and quantitative displays. Bar charts are for categorical; histograms are for quantitative.
- When comparing distributions, you MUST use comparative language.
- Remember: the mean is pulled in the direction of the skew!
\end{trapbox}

\textbf{Formula Reference:}
- Sample Mean: $\bar{x} = \frac{\sum x_i}{n}$
- Standard Deviation: $s_x = \sqrt{\frac{\sum (x_i - \bar{x})^2}{n-1}}$
- Outlier boundaries: $Q_1 - 1.5(IQR)$ and $Q_3 + 1.5(IQR)$
- Z-score: $z = \frac{x - \mu}{\sigma}$
% === UNIT 2 ===
\chapter{Unit 2: Exploring Two-Variable Data}

\section{SECTION A - Unit Overview}
\begin{lessonbox}[Unit 2 Vocabulary]
\begin{itemize}
    \item \textbf{Explanatory Variable}: A variable that helps explain or influences changes in a response variable (x-axis).
    \item \textbf{Response Variable}: A variable that measures an outcome of a study (y-axis).
    \item \textbf{Scatterplot}: Shows the relationship between two quantitative variables measured on the same individuals.
    \item \textbf{Correlation ($r$)}: Measures the direction and strength of the linear relationship between two quantitative variables.
    \item \textbf{LSRL (Least Squares Regression Line)}: The line that minimizes the sum of the squared residuals.
    \item \textbf{Residual}: The difference between an observed value of the response variable and the value predicted by the regression line ($y - \hat{y}$).
    \item \textbf{Residual Plot}: A scatterplot of the residuals against the explanatory variable; helps assess model fit.
    \item \textbf{Coefficient of Determination ($r^2$)}: The fraction of the variation in the values of $y$ that is accounted for by the LSRL.
    \item \textbf{Extrapolation}: The use of a regression line for prediction far outside the interval of values of the explanatory variable.
    \item \textbf{Influential Point}: An observation that, if removed, would markedly change the result of the calculation (especially the LSRL).
\end{itemize}
\end{lessonbox}

\textbf{AP Exam Weight}: 5-7\% of the AP Statistics Exam.

\textbf{What you will learn}: You'll move from exploring single variables to analyzing relationships between two quantitative variables. You will interpret scatterplots, calculate correlation, build Least Squares Regression Lines (LSRL) to make predictions, analyze residuals to check model validity, and handle non-linear data.

\section{SECTION B - Core Concept Lessons}

\begin{lessonbox}[Lesson 2.1: Scatterplots and the DOFS Framework]
\textbf{Definition}: A scatterplot graphs two quantitative variables. When interpreting a scatterplot, always use DOFS: Direction, Outliers, Form, and Strength.

\textbf{Explanation}: 
\textbf{D}irection: Positive (as x increases, y increases) or Negative (as x increases, y decreases).
\textbf{O}utliers: Data points that fall outside the overall pattern.
\textbf{F}orm: Linear (straight line pattern) or Non-linear (curved).
\textbf{S}trength: How closely the points follow the form (Weak, Moderate, Strong).

\textbf{Context}: If plotting height (x) vs weight (y) for adults, you would likely see a positive, strong, linear relationship. Taller people tend to weigh more.
\end{lessonbox}

\begin{lessonbox}[Lesson 2.2: Correlation Coefficient ($r$)]
\textbf{Definition \& Formula}: Correlation ($r$) measures the strength and direction of a \textit{linear} relationship. It is always between -1 and 1. $r = \frac{1}{n-1}\sum(\frac{x_i - \bar{x}}{s_x})(\frac{y_i - \bar{y}}{s_y})$.

\textbf{Explanation}: An $r$ near 0 indicates a weak linear relationship. An $r$ near 1 or -1 indicates a strong linear relationship. Correlation does not imply causation, and it has no units!

\textbf{Context}: The correlation between hours studied and test score might be $r = 0.85$, indicating a strong positive linear relationship.
\end{lessonbox}

\begin{lessonbox}[Lesson 2.3: Least Squares Regression Line (LSRL)]
\textbf{Definition \& Formula}: The LSRL is $\hat{y} = a + bx$. The slope is $b = r(\frac{s_y}{s_x})$, and the y-intercept is $a = \bar{y} - b\bar{x}$.

\textbf{Explanation}: $\hat{y}$ is the \textit{predicted} value. The slope ($b$) tells us the predicted change in $y$ for a 1-unit increase in $x$. The intercept ($a$) is the predicted $y$ when $x=0$.

\textbf{Context}: For an equation predicting salary ($\hat{y}$) from years of experience ($x$): $\hat{y} = 40,000 + 3,000x$. For every additional year of experience, predicted salary increases by \$3,000.
\end{lessonbox}

\begin{lessonbox}[Lesson 2.4: Residuals and $r^2$]
\textbf{Definition \& Formula}: Residual = Actual $y$ - Predicted $y$ ($y - \hat{y}$). $r^2$ is the coefficient of determination.

\textbf{Explanation}: A positive residual means the model underestimated the actual value. A good regression line will have a residual plot with no clear pattern. $r^2$ tells us the percentage of variability in $y$ explained by the linear relationship with $x$.

\textbf{Context}: If predicting test score from study hours yields $r^2 = 0.72$, then 72\% of the variation in test scores is explained by the linear relationship with study hours.
\end{lessonbox}

\section{SECTION C - Worked Examples}

\textbf{Example 1: Interpreting Slope and Intercept}
\textit{Context}: A regression line predicting heating cost ($\hat{y}$, in \$) from average outside temperature ($x$, in $^\circ F$) is $\hat{y} = 300 - 4x$.
\textit{Solution}:
- Slope (-4): For each additional $1^\circ F$ increase in outside temperature, the predicted heating cost decreases by \$4.
- y-intercept (300): On a day where outside temperature is $0^\circ F$, the predicted heating cost is \$300.

\textbf{Example 2: Calculating a Residual}
\textit{Context}: The equation $\hat{y} = 2.5 + 1.2x$ predicts GPA based on study hours per week. A student studies 5 hours and has a GPA of 3.9. Calculate their residual.
\textit{Solution}:
1. Find predicted GPA: $\hat{y} = 2.5 + 1.2(5) = 2.5 + 6.0 = 8.5$. Wait, GPA of 8.5 is invalid. Let's use $\hat{y} = 2.5 + 0.2(5) = 3.5$.
(Recalculating with better numbers: $\hat{y} = 2.5 + 0.2x$. 5 hours $\rightarrow \hat{y} = 2.5 + 0.2(5) = 3.5$.)
2. Actual $y = 3.9$.
3. Residual = $y - \hat{y} = 3.9 - 3.5 = 0.4$.
4. Interpretation: The student's GPA is 0.4 points higher than predicted.

\textbf{Example 3: Interpreting $r^2$}
\textit{Context}: For the relationship between car weight and MPG, $r = -0.90$.
\textit{Solution}:
1. Calculate $r^2$: $(-0.90)^2 = 0.81$.
2. Interpretation: 81\% of the variability in MPG is accounted for by the linear regression on car weight.

\textbf{Example 4: Reading Computer Output}
\textit{Context}: Regression output shows Constant Coef = 15.2, and "Height" Coef = 0.8.
\textit{Solution}:
1. Identify the intercept (Constant) = 15.2.
2. Identify the slope (variable coeff) = 0.8.
3. Equation: $\hat{y} = 15.2 + 0.8(\text{Height})$.

\section{SECTION D - TI-84 CE Keystroke Playbook}
\begin{calculatorbox}[TI-84 Procedures for Unit 2]
\textbf{1. Calculating LSRL and Correlation (LinReg)}
\begin{itemize}
    \item Enter x in L1 and y in L2.
    \item Ensure Diagnostics are ON: \texttt{MODE} $\rightarrow$ \texttt{StatDiagnostics: ON}.
    \item \texttt{STAT} $\rightarrow$ \texttt{CALC} $\rightarrow$ \texttt{8:LinReg(a+bx)}.
    \item Specify Xlist: L1, Ylist: L2. Calculate to get a, b, $r^2$, and $r$.
\end{itemize}

\textbf{2. Creating a Scatterplot}
\begin{itemize}
    \item \texttt{2ND} $\rightarrow$ \texttt{Y=} (STAT PLOT) $\rightarrow$ \texttt{Plot1}.
    \item Type: 1st icon (scatterplot). Xlist: L1, Ylist: L2.
    \item \texttt{ZOOM} $\rightarrow$ \texttt{9:ZoomStat}.
\end{itemize}

\textbf{3. Creating a Residual Plot}
\begin{itemize}
    \item After running LinReg, residuals are auto-saved to RESID.
    \item \texttt{STAT PLOT} $\rightarrow$ \texttt{Plot1}. Xlist: L1, Ylist: \texttt{2ND} $\rightarrow$ \texttt{STAT} (LIST) $\rightarrow$ select \texttt{RESID}.
    \item \texttt{ZOOM} $\rightarrow$ \texttt{9:ZoomStat}.
    \item Look for a random scatter to confirm a linear model is appropriate!
\end{itemize}
\end{calculatorbox}

\section{SECTION E - Score-4 FRQ Templates}

\begin{frqrubricbox}[FRQ Template 1: Interpreting Slope and Y-intercept]
\textbf{Prompt Type:} "Interpret the slope / y-intercept of the LSRL in context."

\textbf{Template for Slope:}
"For each additional [1 unit of x-variable], the \textbf{predicted} [y-variable] [increases/decreases] by [slope value]."

\textbf{Template for Y-intercept:}
"When the [x-variable] is 0 [units], the \textbf{predicted} [y-variable] is [y-intercept value]."

\textbf{Grader Rubric (4 points):}
\begin{itemize}
    \item Must explicitly state the word "predicted", "estimated", or "expected". (Omitting this implies a deterministic relationship and loses points!)
    \item Must include the correct units for both x and y variables.
    \item Must mention the direction (increase/decrease) correctly corresponding to the sign of the slope.
\end{itemize}
\end{frqrubricbox}

\begin{frqrubricbox}[FRQ Template 2: Interpreting a Residual Plot]
\textbf{Prompt Type:} "Is a linear model appropriate for this data? Justify using the residual plot."

\textbf{Template:}
"Yes/No, a linear model is [appropriate / not appropriate] because the residual plot shows [a random scatter of points / a clear curved pattern]."

\textbf{Grader Rubric (4 points):}
\begin{itemize}
    \item Must state clear yes/no conclusion.
    \item Must refer specifically to the \textit{residual plot}, not the scatterplot.
    \item Must describe the pattern (or lack thereof) in the residual plot.
    \item Note: Random scatter = linear model is good. Pattern (like a U-shape) = linear model is bad.
\end{itemize}
\end{frqrubricbox}

\section{SECTION F - Practice Problem Set}
\begin{workbookbox}[Unit 2 Practice Problems]
\textbf{Multiple Choice}
1. Which is true about the correlation coefficient $r$?
(A) It has units (B) It changes if x and y are swapped (C) It measures linear strength (D) $r=0$ means no relationship exists
\textit{Solution: (C). It measures linear strength. A curve could have $r=0$ but strong non-linear relation.}

2. Residual is calculated by:
(A) Predicted - Actual (B) Actual - Predicted (C) Actual / Predicted (D) $x - \bar{x}$
\textit{Solution: (B). $y - \hat{y}$.}

3. If $r = 0.8$, what is $r^2$?
(A) 0.64 (B) 0.8 (C) 1.6 (D) 0.16
\textit{Solution: (A).}

4. Extrapolation occurs when:
(A) You predict inside the domain (B) You predict far outside the domain of x-values (C) The residual is large (D) The correlation is low
\textit{Solution: (B).}

5. A residual plot shows a clear U-shape. This means:
(A) Linear model is appropriate (B) Linear model is NOT appropriate (C) Correlation is 1 (D) $r^2$ is 100\%
\textit{Solution: (B). A pattern in residuals indicates a non-linear relationship.}

6. If slope $b = -2.5$, what must be true about $r$?
(A) $r > 0$ (B) $r < 0$ (C) $r = -2.5$ (D) $r = 1$
\textit{Solution: (B). Slope and correlation always have the same sign.}

7. An influential point is removed. The correlation goes from 0.8 to 0.4. The point was:
(A) Following the pattern (B) Breaking the pattern (C) Exactly on the mean (D) A residual of 0
\textit{Solution: (A). It was artificially inflating the correlation by being far out in the x-direction but aligning with the trend.}

8. Equation: $\hat{y} = 10 + 2x$. If $x=5$, predicted $y$ is:
(A) 10 (B) 20 (C) 15 (D) 0
\textit{Solution: (B). $10 + 2(5) = 20$.}

9. Which transformation helps linearize exponential growth data?
(A) $(x, y^2)$ (B) $(x, \log y)$ (C) $(\log x, y)$ (D) $(\sqrt{x}, y)$
\textit{Solution: (B). Logging the y-variable linearizes exponential data.}

10. The point $(\bar{x}, \bar{y})$ always:
(A) Is an outlier (B) Has a residual of 0 (C) Lies exactly on the LSRL (D) Has maximum leverage
\textit{Solution: (C). The LSRL always passes through the point of averages.}

\textbf{Free Response}
11. Given $\bar{x}=10, s_x=2, \bar{y}=50, s_y=10, r=0.8$. Find the LSRL.
\textit{Solution: $b = r(s_y/s_x) = 0.8(10/2) = 4$. $a = \bar{y} - b\bar{x} = 50 - 4(10) = 10$. LSRL: $\hat{y} = 10 + 4x$.}

12. Interpret $r^2 = 0.64$ in the context of $x=$ age and $y=$ height.
\textit{Solution: 64\% of the variation in height is explained by the linear relationship with age.}

13. The LSRL is $\hat{\text{weight}} = 20 + 5(\text{age})$. A 10-year-old weighs 80 lbs. Find the residual.
\textit{Solution: $\hat{y} = 20 + 5(10) = 70$. Residual = $y - \hat{y} = 80 - 70 = 10$ lbs.}

14. Why is extrapolation dangerous?
\textit{Solution: Because we have no evidence that the linear trend observed in the collected data continues outside that domain.}

15. Describe DOFS for a scatterplot of shoe size vs GPA.
\textit{Solution: Direction: None. Outliers: None expected. Form: No distinct form (random scatter). Strength: Very weak. There is no logical relationship.}
\end{workbookbox}

\section{SECTION G - Unit Summary}
\begin{trapbox}[Common Mistakes]
- Forgetting to write the "hat" on $\hat{y}$ in the regression equation!
- Saying a 1-unit increase in x "causes" an increase in y. Use "is associated with a predicted increase".
- Assuming high correlation means a linear model is appropriate. Always check the residual plot!
\end{trapbox}

\textbf{Formula Reference:}
- LSRL: $\hat{y} = a + bx$
- Slope: $b = r \frac{s_y}{s_x}$
- Intercept: $a = \bar{y} - b\bar{x}$
- Residual: $y - \hat{y}$
% === UNIT 3 ===
\chapter{Unit 3: Collecting Data}

\section{SECTION A - Unit Overview}
\begin{lessonbox}[Unit 3 Vocabulary]
\begin{itemize}
    \item \textbf{Population}: The entire group of individuals we want information about.
    \item \textbf{Sample}: A subset of individuals in the population from which we actually collect data.
    \item \textbf{Observational Study}: Observes individuals and measures variables of interest but does not attempt to influence the responses.
    \item \textbf{Experiment}: Deliberately imposes some treatment on individuals to measure their responses.
    \item \textbf{Confounding}: Occurs when two variables are associated in such a way that their effects on a response variable cannot be distinguished from each other.
    \item \textbf{Simple Random Sample (SRS)}: A sample of size $n$ chosen in such a way that every group of $n$ individuals has an equal chance to be selected.
    \item \textbf{Stratified Random Sample}: Classifies the population into groups of similar individuals (strata), then chooses a separate SRS in each stratum.
    \item \textbf{Cluster Sample}: Divides population into groups (clusters), randomly selects some clusters, and surveys EVERY individual in the selected clusters.
    \item \textbf{Undercoverage Bias}: Occurs when some members of the population cannot be chosen in a sample.
    \item \textbf{Nonresponse Bias}: Occurs when an individual chosen for the sample can't be contacted or refuses to participate.
    \item \textbf{Placebo Effect}: A fake treatment that works because subjects believe it will work.
    \item \textbf{Double-Blind}: Neither the subjects nor those who interact with them and measure the response variable know which treatment a subject received.
\end{itemize}
\end{lessonbox}

\textbf{AP Exam Weight}: 12-15\% of the AP Statistics Exam.

\textbf{What you will learn}: You will learn how to design good surveys and experiments. You'll distinguish between observing data and manipulating it, learn various sampling techniques to avoid bias, and master the four principles of experimental design to establish causation.

\section{SECTION B - Core Concept Lessons}

\begin{lessonbox}[Lesson 3.1: Observational Studies vs. Experiments]
\textbf{Definition}: An \textbf{observational study} gathers data without interfering. An \textbf{experiment} actively imposes a treatment to observe the response.

\textbf{Explanation}: Only a well-designed experiment can establish a cause-and-effect relationship. Observational studies are plagued by confounding variables.

\textbf{Context}: If you want to know if tutoring improves test scores, comparing students who chose to get tutoring with those who didn't is an observational study. Motivation is a confounding variable! To do an experiment, you must randomly assign half the students to receive tutoring.
\end{lessonbox}

\begin{lessonbox}[Lesson 3.2: Sampling Methods]
\textbf{Definition}: 
\textbf{SRS}: Names in a hat. 
\textbf{Stratified}: Split by a trait (e.g., grade level), take some from ALL groups.
\textbf{Cluster}: Split geographically (e.g., homerooms), take ALL from some groups.

\textbf{Explanation}: Stratified sampling guarantees representation from each stratum and reduces variability. Cluster sampling is often done for convenience and cost-effectiveness.

\textbf{Context}: Surveying a high school:
- Stratified: Pick 50 freshmen, 50 sophomores, 50 juniors, 50 seniors.
- Cluster: Randomly pick 5 homerooms, survey every student in those 5 rooms.
\end{lessonbox}

\begin{lessonbox}[Lesson 3.3: Types of Bias]
\textbf{Definition}: Bias is the systematic favoring of certain outcomes. 
\textbf{Voluntary Response}: People choose to respond (usually strong negative opinions).
\textbf{Undercoverage}: Part of population is left out.
\textbf{Nonresponse}: Chosen people refuse.
\textbf{Response}: People lie or questions are worded poorly.

\textbf{Context}: A radio show asks listeners to call in about a new tax. Only angry taxpayers call. This is Voluntary Response Bias. If a survey is done via landline phones, young people are omitted. This is Undercoverage.
\end{lessonbox}

\begin{lessonbox}[Lesson 3.4: Principles of Experimental Design]
\textbf{Definition}: 
1. \textbf{Comparison}: Compare two or more treatments (e.g., control vs treatment).
2. \textbf{Random Assignment}: Use chance to assign treatments to balance out lurking variables.
3. \textbf{Control}: Keep other variables constant for all groups.
4. \textbf{Replication}: Use enough subjects so differences aren't just by chance.

\textbf{Context}: Testing a new fertilizer: Use the same soil and water (Control), randomly assign plants to Fertilizer A or B (Random Assignment), use 100 plants not just 2 (Replication), and compare growth (Comparison).
\end{lessonbox}

\section{SECTION C - Worked Examples}

\textbf{Example 1: Identifying Confounding Variables}
\textit{Context}: A study finds that cities with more churches also have higher crime rates. Can we conclude churches cause crime?
\textit{Solution}:
No. This is an observational study. A confounding variable is population size. Larger cities have both more churches and more crime. We cannot distinguish the effect of churches from the effect of population size.

\textbf{Example 2: How to perform an SRS}
\textit{Context}: You have a list of 800 students. How do you select an SRS of 20?
\textit{Solution}:
1. Label each student with a unique number from 001 to 800.
2. Use a random number generator to select 20 unique numbers between 001 and 800.
3. The students corresponding to these 20 numbers form the sample.

\textbf{Example 3: Stratified vs Cluster}
\textit{Context}: A farmer wants to test the quality of apples in his orchard of 10 rows. 
\textit{Solution}:
- Stratified: He randomly selects 3 apples from every single row (Total = 30 apples). Good if rows get different sunlight.
- Cluster: He randomly selects 1 entire row and checks all 30 apples in that row. Good for saving time walking around.

\textbf{Example 4: Randomized Block Design}
\textit{Context}: Testing a new sunscreen. Men and women sweat differently. How to design?
\textit{Solution}:
1. Block by gender: Separate subjects into a block of men and a block of women.
2. Within the men's block, randomly assign half to the new sunscreen, half to the old.
3. Within the women's block, randomly assign half to the new, half to the old.
4. Compare skin damage within men, and separately within women. This removes variation due to gender.

\section{SECTION D - TI-84 CE Keystroke Playbook}
\begin{calculatorbox}[TI-84 Procedures for Unit 3]
\textbf{1. Setting the Random Seed (Optional but good practice)}
\begin{itemize}
    \item Type any unique number (like your student ID).
    \item Press \texttt{STO $\rightarrow$}.
    \item Press \texttt{MATH} $\rightarrow$ \texttt{PROB} $\rightarrow$ \texttt{1:rand}. Hit \texttt{ENTER}.
\end{itemize}

\textbf{2. Generating Random Integers for an SRS}
\begin{itemize}
    \item \texttt{MATH} $\rightarrow$ \texttt{PROB} $\rightarrow$ \texttt{5:randInt(}
    \item Syntax: \texttt{randInt(lower, upper, n)} where n is how many numbers you want.
    \item Example: \texttt{randInt(1, 800, 20)} generates 20 numbers between 1 and 800.
    \item \textit{Note: You must manually discard any repeat numbers if sampling without replacement!}
\end{itemize}

\textbf{3. Generating Random Integers without Repeats (Newer TI-84 CE)}
\begin{itemize}
    \item \texttt{MATH} $\rightarrow$ \texttt{PROB} $\rightarrow$ \texttt{8:randIntNoRep(}
    \item Syntax: \texttt{randIntNoRep(lower, upper)}
    \item Generates a list of all integers in the range in a random scrambled order. Just take the first $n$ numbers.
\end{itemize}
\end{calculatorbox}

\section{SECTION E - Score-4 FRQ Templates}

\begin{frqrubricbox}[FRQ Template 1: Describing an SRS]
\textbf{Prompt Type:} "Describe how you would select a Simple Random Sample of size $n$ from a population of size $N$."

\textbf{Template:}
"First, I will assign a unique integer from 1 to [N] to each [individual in population]. Using a random number generator, I will generate [n] \textbf{unique} random integers from 1 to [N]. I will ignore any repeated numbers. The [n] individuals whose labels correspond to the generated numbers will be selected for the sample."

\textbf{Grader Rubric (4 points):}
\begin{itemize}
    \item \textbf{Labeling}: Must state that numbers are assigned to individuals.
    \item \textbf{Range}: Must specify the range of numbers used.
    \item \textbf{Selection}: Must state how numbers are chosen (RNG or hat).
    \item \textbf{Repeats}: Must explicitly state that repeated numbers are skipped/ignored! (Crucial point).
\end{itemize}
\end{frqrubricbox}

\begin{frqrubricbox}[FRQ Template 2: Explaining Bias]
\textbf{Prompt Type:} "Explain why this sampling method is biased and the direction of the bias."

\textbf{Template:}
"This sampling method suffers from [Type of Bias] because [explain how the sample doesn't represent the population]. As a result, the sample will likely have a disproportionate number of [specific group]. Because this group tends to [explain their specific trait related to the response variable], the survey will likely [overestimate / underestimate] the true [parameter being measured]."

\textbf{Grader Rubric (4 points):}
\begin{itemize}
    \item Identify the specific bias or flaw in the design.
    \item Connect the flaw to the specific context of the problem.
    \item State whether the estimate will be an overestimation or underestimation.
    \item Explain \textit{why} the flaw leads to that specific direction of bias.
\end{itemize}
\end{frqrubricbox}

\section{SECTION F - Practice Problem Set}
\begin{workbookbox}[Unit 3 Practice Problems]
\textbf{Multiple Choice}
1. To establish cause-and-effect, you must use:
(A) An observational study (B) A census (C) A well-designed experiment (D) A stratified sample
\textit{Solution: (C). Only experiments can show causation.}

2. Blocking in an experiment is similar to what in sampling?
(A) SRS (B) Stratified sampling (C) Cluster sampling (D) Systematic sampling
\textit{Solution: (B). Both group similar individuals together to reduce variation.}

3. A survey on a website asks visitors to vote on a policy. This is:
(A) SRS (B) Voluntary response (C) Cluster (D) Stratified
\textit{Solution: (B). Visitors self-select to participate.}

4. Purpose of random assignment in an experiment:
(A) Create roughly equivalent groups (B) Eliminate the placebo effect (C) Ensure sample represents population (D) Reduce sample size
\textit{Solution: (A). It balances out confounding variables between treatment groups.}

5. A fake pill used in medical trials is called:
(A) A block (B) A treatment (C) A placebo (D) A stratum
\textit{Solution: (C).}

6. A pollster randomly calls landline phones during the day. What bias is most likely?
(A) Nonresponse bias (B) Undercoverage bias (C) Response bias (D) Both A and B
\textit{Solution: (D). Undercoverage (no cell phones) and nonresponse (people not home).}

7. We group students by grade level and take a random sample of 10 from each grade. This is:
(A) Cluster (B) Stratified (C) SRS (D) Systematic
\textit{Solution: (B). Some from all groups.}

8. Which principle of design reduces variation by applying treatments to many subjects?
(A) Control (B) Random Assignment (C) Replication (D) Comparison
\textit{Solution: (C).}

9. Double-blind means:
(A) Subjects don't know the treatment (B) Evaluators don't know the treatment (C) Both A and B (D) Two placebos are used
\textit{Solution: (C).}

10. You survey every 10th person walking into a store. This is:
(A) SRS (B) Stratified (C) Systematic (D) Cluster
\textit{Solution: (C).}

\textbf{Free Response}
11. Explain the difference between stratified and cluster sampling.
\textit{Solution: Stratified divides population into groups and takes a random sample from EVERY group. Cluster divides population into groups, randomly selects SOME groups, and censuses EVERY individual in those selected groups.}

12. A gym wants to test a new protein powder. Explain a confounding variable if they just compare members who buy it vs those who don't.
\textit{Solution: Members who buy the powder might also work out more frequently or have better overall diets. We wouldn't know if muscle gain is from the powder or the extra workouts.}

13. Design a completely randomized design for testing 2 drugs on 100 patients.
\textit{Solution: Assign patients numbers 1-100. Use a RNG to select 50 unique numbers. These 50 patients receive Drug A. The remaining 50 receive Drug B. Compare the recovery rates.}

14. Why is blinding important in an experiment?
\textit{Solution: It prevents the placebo effect in subjects and prevents measurement bias from the researchers who are evaluating the results.}

15. What is the main benefit of blocking?
\textit{Solution: Blocking isolates the variation caused by the blocking variable, making it easier to detect a difference between the treatments.}
\end{workbookbox}

\section{SECTION G - Unit Summary}
\begin{trapbox}[Common Mistakes]
- Using "random assignment" and "random sampling" interchangeably. Random sampling is how you get your subjects (generalizability). Random assignment is how you assign treatments (causation).
- Forgetting to say "ignore repeats" when describing an SRS.
- Confusing stratified sampling (some from all) with cluster sampling (all from some).
\end{trapbox}

\textbf{Key Takeaways:}
- Observational Study $\rightarrow$ Association only.
- Experiment $\rightarrow$ Causation.
- Random Selection $\rightarrow$ Allows generalization to population.
- Random Assignment $\rightarrow$ Allows cause-and-effect conclusion.
% === UNIT 4 ===
\chapter{Unit 4: Probability and Random Variables}

\section{SECTION A - Unit Overview}
\begin{lessonbox}[Unit 4 Vocabulary]
\begin{itemize}
    \item \textbf{Sample Space ($S$)}: The set of all possible outcomes of a chance process.
    \item \textbf{Event}: Any collection of outcomes from some chance process.
    \item \textbf{Mutually Exclusive (Disjoint)}: Two events that have no outcomes in common and can never occur together.
    \item \textbf{Independent Events}: Two events are independent if the occurrence of one does not change the probability that the other will happen.
    \item \textbf{Conditional Probability}: The probability that one event happens given that another event is already known to have happened, $P(A | B)$.
    \item \textbf{Random Variable (RV)}: Takes numerical values that describe the outcomes of some chance process.
    \item \textbf{Discrete RV}: Can take a fixed set of possible values with gaps between them (usually counts).
    \item \textbf{Continuous RV}: Can take any value in an interval on the number line (usually measurements).
    \item \textbf{Expected Value (Mean of RV)}: The long-run average outcome of a random variable.
    \item \textbf{Binomial Setting}: Arises when we perform several independent trials of the same chance process and record the number of times that a particular outcome occurs.
    \item \textbf{Geometric Setting}: Arises when we perform independent trials of the same chance process and record the number of trials until a particular outcome occurs.
\end{itemize}
\end{lessonbox}

\textbf{AP Exam Weight}: 10-20\% of the AP Statistics Exam.

\textbf{What you will learn}: You will learn the formal rules of probability, how to evaluate conditional probabilities, and how to use Venn diagrams and tree diagrams. You'll define random variables, calculate expected values, combine distributions, and identify and calculate probabilities for Binomial and Geometric distributions.

\section{SECTION B - Core Concept Lessons}

\begin{lessonbox}[Lesson 4.1: General Probability Rules]
\textbf{Definition \& Formula}: 
\textbf{Complement Rule}: $P(A^c) = 1 - P(A)$.
\textbf{General Addition Rule}: $P(A \cup B) = P(A) + P(B) - P(A \cap B)$.
\textbf{General Multiplication Rule}: $P(A \cap B) = P(A) \times P(B | A)$.

\textbf{Explanation}: The addition rule uses "OR" ($\cup$), meaning at least one event occurs. We subtract the intersection so we don't double count! The multiplication rule uses "AND" ($\cap$).

\textbf{Context}: Probability of being left-handed is 0.10. Probability of having blue eyes is 0.30. If these traits are independent, probability of having BOTH is $0.10 \times 0.30 = 0.03$.
\end{lessonbox}

\begin{lessonbox}[Lesson 4.2: Conditional Probability and Independence]
\textbf{Definition \& Formula}: 
Conditional formula: $P(A | B) = \frac{P(A \cap B)}{P(B)}$.
Independence check: $A$ and $B$ are independent if $P(A | B) = P(A)$.

\textbf{Explanation}: Independence means knowing $B$ happened gives you NO new information about $A$. Do not confuse independent with mutually exclusive. Mutually exclusive events are highly DEPENDENT (if one happens, the other has 0\% chance of happening).

\textbf{Context}: If $P(\text{A on test}) = 0.8$, but $P(\text{A on test} | \text{slept well}) = 0.95$, then getting an A and sleeping well are NOT independent.
\end{lessonbox}

\begin{lessonbox}[Lesson 4.3: Transforming and Combining Random Variables]
\textbf{Definition \& Formula}: 
Linear Transformation: If $Y = a + bX$, then $\mu_Y = a + b\mu_X$ and $\sigma_Y = |b|\sigma_X$.
Combining: If $X$ and $Y$ are independent, $\mu_{X+Y} = \mu_X + \mu_Y$ and $\sigma^2_{X+Y} = \sigma^2_X + \sigma^2_Y$.

\textbf{Explanation}: Means always add/subtract directly. Variances ALWAYS add, even if you are subtracting the random variables! You CANNOT add standard deviations directly.

\textbf{Context}: Packing a box. Weight of box is $X$, weight of item is $Y$. Total weight $= X+Y$. The variance of the total weight is the variance of the box plus the variance of the item.
\end{lessonbox}

\begin{lessonbox}[Lesson 4.4: The Binomial Distribution (BINS)]
\textbf{Definition \& Formula}: 
Check \textbf{BINS}:
\textbf{B}inary (Success/Failure)
\textbf{I}ndependent trials
\textbf{N}umber of trials is fixed ($n$)
\textbf{S}ame probability of success ($p$).
Formula: $P(X=k) = \binom{n}{k} p^k (1-p)^{n-k}$.

\textbf{Explanation}: The binomial distribution counts the number of successes in a fixed number of trials. Mean $\mu = np$, standard deviation $\sigma = \sqrt{np(1-p)}$.

\textbf{Context}: Flipping a coin 10 times and counting heads. $n=10, p=0.5$. What is the probability of exactly 6 heads?
\end{lessonbox}

\section{SECTION C - Worked Examples}

\textbf{Example 1: Using the Addition Rule}
\textit{Context}: At a school, 40\% of students take French, 30\% take Spanish, and 10\% take both. Find the probability a random student takes French OR Spanish.
\textit{Solution}:
$P(F \cup S) = P(F) + P(S) - P(F \cap S)$
$P(F \cup S) = 0.40 + 0.30 - 0.10 = 0.60$.

\textbf{Example 2: Checking for Independence in a Table}
\textit{Context}: 
Event $M$ = Male (50 out of 100). Event $S$ = Plays sports (40 out of 100). 25 are Males who play sports. Are $M$ and $S$ independent?
\textit{Solution}:
Check if $P(S | M) = P(S)$.
$P(S) = 40/100 = 0.40$.
$P(S | M) = \frac{P(S \cap M)}{P(M)} = \frac{25/100}{50/100} = \frac{25}{50} = 0.50$.
Since $0.50 \neq 0.40$, they are NOT independent.

\textbf{Example 3: Expected Value}
\textit{Context}: A lottery ticket costs \$5. It has a 10\% chance of winning \$20, a 1\% chance of winning \$100, and otherwise pays \$0. What is the expected profit?
\textit{Solution}:
Net outcomes: Win \$20 $\rightarrow$ Profit \$15. Win \$100 $\rightarrow$ Profit \$95. Lose $\rightarrow$ Profit -\$5.
$P(15) = 0.10$, $P(95) = 0.01$, $P(-5) = 0.89$.
$E(X) = (15)(0.10) + (95)(0.01) + (-5)(0.89) = 1.50 + 0.95 - 4.45 = -2.00$.
Expected profit is -\$2.00.

\textbf{Example 4: Binomial Probability}
\textit{Context}: A basketball player makes 80\% of his free throws. He shoots 5 independent free throws. What is the probability he makes exactly 3?
\textit{Solution}:
$n=5, p=0.8, k=3$.
$P(X=3) = \binom{5}{3} (0.8)^3 (0.2)^2 = (10)(0.512)(0.04) = 0.2048$.

\section{SECTION D - TI-84 CE Keystroke Playbook}
\begin{calculatorbox}[TI-84 Procedures for Unit 4]
\textbf{1. binompdf (Exact binomial probability)}
\begin{itemize}
    \item Finds $P(X = k)$.
    \item \texttt{2ND} $\rightarrow$ \texttt{VARS} (DISTR) $\rightarrow$ \texttt{A:binompdf(}
    \item Syntax: \texttt{binompdf(trials, p, x value)}
    \item Example: \texttt{binompdf(5, 0.8, 3)} yields 0.2048.
\end{itemize}

\textbf{2. binomcdf (Cumulative binomial probability)}
\begin{itemize}
    \item Finds $P(X \leq k)$.
    \item \texttt{2ND} $\rightarrow$ \texttt{VARS} (DISTR) $\rightarrow$ \texttt{B:binomcdf(}
    \item Syntax: \texttt{binomcdf(trials, p, x value)}
    \item Example: To find $P(X \geq 3)$, do $1 - \texttt{binomcdf(5, 0.8, 2)}$.
\end{itemize}

\textbf{3. geometpdf / geometcdf}
\begin{itemize}
    \item \texttt{2ND} $\rightarrow$ \texttt{VARS} (DISTR) $\rightarrow$ \texttt{E:geometpdf(} or \texttt{F:geometcdf(}
    \item Syntax: \texttt{geometpdf(p, x value)}
    \item Used for Geometric distributions (counting trials UNTIL a success).
\end{itemize}
\end{calculatorbox}

\section{SECTION E - Score-4 FRQ Templates}

\begin{frqrubricbox}[FRQ Template 1: Identifying a Binomial Distribution]
\textbf{Prompt Type:} "Explain why the random variable $X$ follows a binomial distribution."

\textbf{Template:}
"$X$ is a binomial random variable because it satisfies the four BINS criteria:
1. \textbf{B}inary: Each trial results in a success ([define success]) or failure ([define failure]).
2. \textbf{I}ndependent: The outcome of one trial does not affect the outcome of another.
3. \textbf{N}umber fixed: There is a fixed number of trials, $n = $ [number].
4. \textbf{S}ame probability: The probability of success, $p = $ [probability], is the same for each trial."

\textbf{Grader Rubric (4 points):}
\begin{itemize}
    \item Must explicitly list all 4 criteria.
    \item Must define success/failure in context.
    \item Must state the values of $n$ and $p$.
    \item Just listing "BINS" without context will earn a maximum of 1 point!
\end{itemize}
\end{frqrubricbox}

\begin{frqrubricbox}[FRQ Template 2: Expected Value Interpretation]
\textbf{Prompt Type:} "Calculate and interpret the expected value of $X$."

\textbf{Template:}
"The expected value of $X$ is [calculated value]. If we were to repeat this chance process \textbf{many, many times}, the \textbf{long-run average} [context of random variable] would be approximately [calculated value]."

\textbf{Grader Rubric (4 points):}
\begin{itemize}
    \item Correct calculation with work shown.
    \item Must use phrases indicating long-run behavior ("many times", "in the long run").
    \item Must use the word "average" or "mean".
    \item Must include context and units.
\end{itemize}
\end{frqrubricbox}

\section{SECTION F - Practice Problem Set}
\begin{workbookbox}[Unit 4 Practice Problems]
\textbf{Multiple Choice}
1. If $A$ and $B$ are mutually exclusive, $P(A \cap B)$ is:
(A) $P(A) \times P(B)$ (B) 0 (C) 1 (D) $P(A) + P(B)$
\textit{Solution: (B). They cannot occur together.}

2. $P(A)=0.3, P(B)=0.5$. If independent, $P(A \cup B) =$ ?
(A) 0.8 (B) 0.15 (C) 0.65 (D) 0.5
\textit{Solution: (C). $0.3 + 0.5 - (0.3)(0.5) = 0.8 - 0.15 = 0.65$.}

3. You roll a die until you get a 6. This is a:
(A) Binomial setting (B) Geometric setting (C) Normal setting (D) Continuous setting
\textit{Solution: (B).}

4. For a binomial variable with $n=10, p=0.4$, the mean is:
(A) 4 (B) 6 (C) 10 (D) 0.4
\textit{Solution: (A). $\mu = np = 10(0.4) = 4$.}

5. Variance of $X$ is 10. Variance of $Y$ is 5. If independent, variance of $X-Y$ is:
(A) 5 (B) 15 (C) 50 (D) 2
\textit{Solution: (B). Variances ALWAYS add! $10+5=15$.}

6. $P(\text{Rain})=0.2$. $P(\text{Wind})=0.3$. $P(\text{Rain} \cap \text{Wind})=0.1$. Find $P(\text{Rain} | \text{Wind})$.
(A) 0.5 (B) 0.2 (C) 0.33 (D) 0.67
\textit{Solution: (C). $0.1 / 0.3 = 1/3 \approx 0.33$.}

7. A random variable $X$ has mean 5. If $Y = 3X + 2$, what is mean of $Y$?
(A) 5 (B) 15 (C) 17 (D) 2
\textit{Solution: (C). $3(5)+2 = 17$.}

8. If standard deviation of $X$ is 4, standard deviation of $Y = 3X + 2$ is:
(A) 12 (B) 14 (C) 16 (D) 38
\textit{Solution: (A). Adding 2 does not affect spread. Multiplying by 3 scales standard dev by 3. $3 \times 4 = 12$.}

9. Which is true for a continuous random variable?
(A) $P(X = c) > 0$ (B) $P(X = c) = 0$ (C) There are gaps (D) Can only be integers
\textit{Solution: (B). Area under a curve at a single point is 0.}

10. A discrete random variable must have probabilities that sum to:
(A) 100 (B) 1 (C) 0 (D) $\mu$
\textit{Solution: (B).}

\textbf{Free Response}
11. $P(A) = 0.4, P(B) = 0.5, P(A \cup B) = 0.7$. Are $A$ and $B$ independent?
\textit{Solution: $P(A \cup B) = P(A) + P(B) - P(A \cap B) \Rightarrow 0.7 = 0.4 + 0.5 - P(A \cap B) \Rightarrow P(A \cap B) = 0.2$. If independent, $P(A \cap B) = P(A)P(B) = (0.4)(0.5) = 0.2$. Yes, they are independent.}

12. A test has 5 multiple choice questions with 4 options each. If you guess all 5, what is the probability you get exactly 2 correct?
\textit{Solution: Binomial. $n=5, p=0.25, k=2$. $P(X=2) = \binom{5}{2}(0.25)^2(0.75)^3 \approx 0.2637$.}

13. Using the previous question, what is the expected number of correct answers?
\textit{Solution: $\mu = np = 5(0.25) = 1.25$ correct answers.}

14. You draw cards from a deck without replacement until you get an Ace. Is this Geometric? Why or why not?
\textit{Solution: No. Because it's without replacement, the probability changes each time. The trials are not independent.}

15. $X$ has mean 10, std dev 2. $Y$ has mean 15, std dev 3. Independent. Find mean and std dev of $X+Y$.
\textit{Solution: Mean $= 10+15 = 25$. Variance $= 2^2 + 3^2 = 4 + 9 = 13$. Std Dev $= \sqrt{13} \approx 3.61$.}
\end{workbookbox}

\section{SECTION G - Unit Summary}
\begin{trapbox}[Common Mistakes]
- Adding standard deviations instead of adding VARIANCES. Always square them first!
- Assuming "mutually exclusive" means "independent". It's the exact opposite!
- Forgetting to show work for binomial probabilities. Just writing "binompdf(10, 0.5, 3)" will not get full credit on the FRQ. You must write out the formula $\binom{10}{3}(0.5)^3(0.5)^7$.
\end{trapbox}

\textbf{Formula Reference:}
- Conditional Prob: $P(A | B) = \frac{P(A \cap B)}{P(B)}$
- Binomial Prob: $P(X=k) = \binom{n}{k} p^k (1-p)^{n-k}$
- Mean of Binomial: $\mu = np$
- Std Dev of Binomial: $\sigma = \sqrt{np(1-p)}$
- Variances sum: $\sigma^2_{X \pm Y} = \sigma^2_X + \sigma^2_Y$
